\chapter{Background}
\label{ch:background}

The middleware presented here moves one data structure, the key-value cache,
between two GPU servers, so its design depends on what a GPU is, on how the
cache arises inside a language model, and on the software that reads it. The
background follows the order in which the later chapters use it: the GPU and
its generations, the arithmetic and number
formats of a model, the transformer and its attention mechanism, the kernels
that compute attention on each generation of hardware, and finally inference
as a systems problem, where the cache and the idea of splitting the work
across machines appear.

\section{Graphics Processors and Their Generations}
\label{ch:gpu}
\label{ch:gpugen}

The processor that runs a language model was designed to draw triangles, and
the properties it kept from that purpose explain both why it is fast and why
the cache it produces is so large. Drawing a three-dimensional scene repeats
one small computation a very large number of times, since every corner of
every triangle is moved by the same matrix arithmetic and every pixel is
colored by the same rule, and no pixel needs the result of another. Hardware
built for this work has many simple arithmetic units side by side and a great
deal of memory bandwidth to feed them. The step that mattered for everything
afterwards was making that hardware programmable: a vertex engine that ran a
small user-written program on every corner \citep{lindholm2001} turned the
graphics card into a machine that applies an arbitrary short program to a
large array of independent inputs.

Researchers soon used such a machine for more than pictures, at first by
disguising linear algebra, simulation and database operations as drawing
\citep{owens2007}, which worked but forced every problem into the vocabulary
of graphics. CUDA removed the disguise by letting a programmer write the
per-element program as an ordinary function, launch it over an array of
threads, and manage GPU memory directly \citep{nickolls2008}. The workload
that made this matter was the neural network, because training a deep network
is dominated by the same multiply-and-add over large arrays, and a network
trained on GPUs won a major image-recognition benchmark by a wide margin
\citep{krizhevsky2017}. Since then the size of models has been set by how
much arithmetic and memory GPUs can supply. Two properties follow for this
thesis. The memory of a GPU is separate from that of the host, fast and
limited, so the cache lives there and moving it to another machine means
leaving that memory (\cref{ch:compat}). In addition, the hardware changes
with every generation, so software is compiled for a generation.

NVIDIA identifies each architecture by a compute capability number, and a
kernel compiled for one number may use instructions that an older number does
not have \citep{cuda_gpus}. \Cref{tab:cc} lists the GPUs in the testbed with
two data-center parts for comparison.

\begin{table}[htbp]
  \centering
  \caption[Compute capability of the GPUs in the testbed]{Compute capability of the GPUs in the testbed, with two data-center parts for reference. The Volta value is not on NVIDIA's current list, which points to a separate list for earlier architectures.}
  \label{tab:cc}
  \small
  \begin{tabular}{L{4.4cm}L{3.2cm}L{2.6cm}L{3.2cm}}
    \toprule
    GPU & Architecture & Compute capability & Where \\
    \midrule
    Tesla V100 & Volta & 7.0 & server 4 \\
    GeForce RTX 3090 & Ampere & 8.6 & server 2 \\
    H100 & Hopper & 9.0 & reference only \\
    B200 & Blackwell, data center & 10.0 & reference only \\
    GeForce RTX 5090, RTX 5070 Ti & Blackwell, consumer & 12.0 & servers 1 and 3 \\
    \bottomrule
  \end{tabular}
\end{table}

Three points from the table are used later. The testbed spans the range,
from the V100, which is the oldest architecture and needs a legacy driver
branch (\cref{ch:bringup}), to the RTX 5090 and 5070 Ti, which are the
newest. Data-center and consumer parts of the newest generation carry
different compute capabilities, so software released for the data-center
part does not necessarily recognize the consumer part, which is the
background to the start-up failure on server 1 (\cref{ch:kernels}). Finally,
consumer and data-center cards differ in more than speed, since features
such as direct access to GPU memory from a network card are reserved for
data-center and workstation products (\cref{ch:compat}). Each generation
since Volta has also added or reworked dedicated matrix-multiply hardware,
called tensor cores, and the instructions that drive it
\citep{patel2025tensorcore}, while the precision of the arithmetic has moved
from 32-bit floating point to 16-bit, and for some uses to 8 bits or fewer
\citep{micikevicius2018,rouhani2023mx}.

\section{Arithmetic and Number Formats}
\label{ch:matrix}

Everything a language model computes is built from a small set of operations
on vectors and matrices, and the size of the cache in later chapters is
simply a count of the numbers that these operations produce. The dot product
of two vectors $a$ and $b$ of the same length $d$ multiplies them position by
position and adds the results:
\mathalt{a dot b equals the sum over i from 1 to d of a i times b i}{%
\begin{equation}
  a \cdot b = \sum_{i=1}^{d} a_i b_i ,
  \label{eq:dot}
\end{equation}}
where $a_i$ and $b_i$ are the entries of the two vectors at position $i$ and
the sum runs over all $d$ positions. The dot product measures how much two
vectors point the same way, being large and positive when they agree, near
zero when they are unrelated, and negative when they oppose, and attention is
built on that reading of it. Multiplying a matrix by a vector takes the dot
product of each row with the vector, so a matrix with $n$ rows and $d$
columns costs $n d$ multiplications and as many additions, none of which
depends on another, which is the shape of work a GPU is built for. A layer
of a neural network multiplies its input by a matrix of learned weights and
applies a simple nonlinear function to each result, and the size of a model
is mostly the number of entries in its weight matrices.

Several times in a model a list of arbitrary scores has to become a list of
positive weights that add up to one, and the softmax function does this. For
a list of $n$ scores $z = (z_1, \dots, z_n)$ it returns the $n$ weights
\mathalt{softmax of z, component i, equals e to the z i, divided by the sum over j from 1 to n
  of e to the z j}{%
\begin{equation}
  \operatorname{softmax}(z)_i = \frac{e^{z_i}}{\sum_{j=1}^{n} e^{z_j}} ,
  \qquad i = 1, \dots, n,
  \label{eq:softmax}
\end{equation}}
where $z_i$ is the score at position $i$, $e$ is the base of the natural
logarithm, and the sum in the denominator runs over all $n$ scores. The
largest score receives the largest weight and the gaps are exaggerated,
so that scores of $(2, 1, 0)$ give weights of about $(0.67, 0.24, 0.09)$.

A number in a model is stored in a fixed number of bits, and the choice sets
both the accuracy of the arithmetic and the size of everything the model
keeps. \Cref{tab:formats} lists the formats that appear in later chapters.

\begin{table}[htbp]
  \centering
  \caption[Number formats and their size for one cache vector of 64 numbers]{Number formats and their size for one cache vector of 64 numbers.}
  \label{tab:formats}
  \small
  \begin{tabular}{L{5.4cm}L{2.6cm}L{4.4cm}}
    \toprule
    Format & Bits per number & Bytes for 64 numbers \\
    \midrule
    32-bit floating point & 32 & 256 \\
    16-bit floating point, two variants & 16 & 128 \\
    8-bit & 8 & 64 \\
    Packed format of \cref{ch:shrink}, with its sketch & 3, plus overhead & 64 \\
    \bottomrule
  \end{tabular}
\end{table}

Models were once trained and run in 32-bit floating point. Training in a mix
of 16-bit and 32-bit arithmetic roughly halves memory and speeds up the
hardware with little loss of accuracy \citep{micikevicius2018}, so 16 bits is
now the usual format for inference, and formats of 8 bits and fewer, with a
shared scale per small group of numbers, go further \citep{rouhani2023mx}.
The model used in the experiments runs in a 16-bit format, so each number in
its cache occupies two bytes.

\section{The Transformer}
\label{ch:transformers}

A language model reads a sequence of tokens and predicts the next one, and
the architecture used for this purpose in every current large model is the
transformer \citep{vaswani2017}. The remaining chapters depend on one part
of it, attention, which is described here in sufficient detail to show what
must be remembered between one token and the next.

\subsection{Tokens and Vectors}

Text is cut into tokens, pieces of words drawn from a fixed vocabulary, and
each token is replaced by a vector of a few thousand numbers. The model is a
stack of identical layers. Each layer takes one vector per position and
returns one vector per position, and the output of the last layer at the last
position is turned into a score for every word in the vocabulary.

\subsection{Attention}

Inside a layer, positions exchange information through attention. Each position computes three vectors from its current vector by three learned
matrices: a query, which expresses what the position is looking for, a key,
which expresses what the position offers, and a value, which is what it hands
over.

For a position with query $q$, the score of an earlier position with key $k$
is their dot product. The scores over all earlier positions are passed
through a softmax to become weights, and the position's result is the
weighted sum of the values. Written for all positions at once,
\mathalt{attention of Q, K and V equals softmax of Q times K transpose, divided by the
  square root of d, all times V}{%
\begin{equation}
  \operatorname{Attention}(Q, K, V) = \operatorname{softmax}\!\left(\frac{Q K^{\top}}{\sqrt{d}}\right) V ,
  \label{eq:attention}
\end{equation}}
where $Q$, $K$ and $V$ are the matrices whose rows are the queries, keys and
values of the positions, $K^{\top}$ is the transpose of $K$, so that the
entry of $Q K^{\top}$ in row $i$ and column $j$ is the dot product of query
$i$ with key $j$, $d$ is the number of entries of a key, and the softmax of
\cref{eq:softmax} is applied to each row. The division by $\sqrt{d}$ keeps
the scores in a range where the softmax behaves well \citep{vaswani2017}.

\subsection{Attention Heads}

One set of queries, keys and values can capture only one kind of relationship, so a
layer runs several in parallel, called heads, each with its own matrices and
shorter vectors, and joins their results. The model used in the experiments has vectors of 64 numbers per head.

Keeping a key and a value for every head, position and layer is expensive, and
two refinements reduce it. In multi-query attention all heads of a layer share
one key and one value per position \citep{shazeer2019mqa}. Grouped-query
attention is the middle ground: the query heads are divided into groups, and
each group shares one key head and one value head \citep{ainslie2023gqa}. The
model used here has 8 key-value heads per layer. What matters for the transfer problem is the count of key-value heads, because keys and values are what must be
kept.

\subsection{The Causal Mask}

A language model is trained to predict each token from the ones before it, so
a position may attend only to itself and earlier positions. The scores for
later positions are replaced by minus infinity before the softmax, which
makes their weights zero. This is the causal mask.

The mask has a consequence that the whole thesis rests on. The vector at a
position, in every layer, depends only on the tokens up to that position.
Adding a token at the end changes nothing that was computed for earlier
positions. Their keys and values can be computed once and kept.

\subsection{The Remainder of a Layer}

After attention, the vector of each position passes through a small two-layer network of the kind described in \cref{ch:matrix}, on its own, with no
exchange between positions. Both steps add their result to their input, and
the vectors are rescaled along the way. None of this involves other
positions, so it adds nothing that must be remembered.

\section{Attention Kernels}
\label{ch:kernels}

Attention is computed by a kernel, a routine compiled for one generation of
GPU, and the kernel matters to a thesis on transport for one reason: it is
the consumer of the cache. It determines the layout and the number format of
the bytes that cross the link, it is compiled separately for each GPU
generation in the testbed, and it was the source of the one software failure
that kept a server out of the cluster. Only what the later chapters rely on
is covered here.

\subsection{The FlashAttention Line of Work}

Attention, as \cref{ch:transformers} defined it, forms a score for every pair
of positions. Computed naively, that is a matrix with one entry per pair,
written to GPU memory and read back. FlashAttention \citep{dao2022fa1}
computes the same result in tiles small enough to stay in fast on-chip
memory, so the full matrix is never stored. The result is exact, and the
saving is in reads and writes of GPU memory, which is the scarce resource
(\cref{sec:history}). FlashAttention-2 \citep{dao2023fa2} reorganized the work
for better parallelism. FlashAttention-3 \citep{shah2024fa3} targets a newer
GPU generation, overlapping data movement with computation and adding a
low-precision mode. FlashAttention-4 \citep{zadouri2026fa4} continues the
line for the most recent hardware.

Two consequences matter here. First, these kernels are exact, so a cache
produced under one of them is valid input to another. That is what allows
prefill on one GPU generation and decode on a different one
(\cref{ch:compat}). Second, each version depends on hardware instructions
that exist only from some generation onward, so which kernel runs is decided
per machine.

\subsection{One Kernel per GPU Generation}

NVIDIA labels each GPU architecture with a compute capability number, and
compiled kernels are specific to it \citep{cuda_gpus}. The testbed spans
three of them (\cref{tab:cc}).

The consumer and data-center parts of the newest generation carry different
numbers. The vendor's list gives 12.0 for the cards of server 1. The distinction matters because software that
checks for ``the newest generation'' by one number will miss the other. The
tensor-core hardware these kernels drive has also changed with each
generation \citep{patel2025tensorcore}.

The difference between generations is also a practical one. On server 1 the
serving engine stopped with an initialization error when the third
FlashAttention version was selected, because one of its modules imported the
attention library, and with it the GPU runtime, before the engine had set up
its GPU context, which the older generation tolerated and the newest did not.
Moving that import into the function that uses it is a one-file patch, kept
in the repository with a container recipe that applies it
\citep{scuffedrdma}, and without it server 1 could not serve at all.

\subsection{Constraints the Kernel Places on the Cache}

The serving engine stores the cache in fixed-size blocks of 16 tokens and
maps each request to a list of blocks \citep{kwon2023vllm}. The attention
kernel reads keys and values straight from those blocks. The connector of \cref{sec:connector} therefore gathers blocks, not tensors, and has to understand
the block layout of the engine version it runs against. It supports three
layouts, and each has a gather and scatter test on both the processor and the
GPU (\cref{sec:tests}).

The number format is the kernel's as well: 16-bit floating point in this
work. Lower-precision formats exist for both weights and cache
\citep{rouhani2023mx,micikevicius2018}, and the engine can store the cache in
8 bits. The quantizer of \cref{ch:shrink} is different in kind. It does not
change what the kernel reads. It packs the bytes for the wire and for disk,
and they are expanded back to 16-bit before the kernel sees them. That
expansion is the one GPU routine written for this work, and its timing is
in \cref{sec:res-decomp}.

\section{LLM Inference as a Systems Problem}
\label{ch:inference}

Running a trained model is called inference. For a language model, inference
is a loop that produces one token at a time, and almost everything that makes
it a systems problem follows from a single data structure that the loop
builds and reuses. That structure comes first, followed by the two phases
of the loop, the metrics visible to users, and the idea of running the two
phases on different machines, which is the point at which the bottleneck
of \cref{sec:history} moves from memory to the link.

\subsection{Autoregressive Generation}

To answer a prompt the model predicts the first token of the answer, appends
it to the sequence, predicts the next, and so on. Done naively, each step
would run the whole sequence through every layer again, and the work to
produce an answer would grow with the square of its length in every layer.

\subsection{The Key-Value Cache}

The causal mask (\cref{ch:transformers}) makes most of that work unnecessary.
The keys and values of earlier positions do not change when a token is added,
so they can be stored the first time they are computed. This store is the
key-value cache, or KV cache. With it, a step computes a query, key and value
for the new position only, and attention reads the stored keys and values of
the earlier ones.

\begin{proposition}[The cache is exact]
\label{prop:cache}
Let a transformer with a causal mask process a sequence of $t$ tokens, and
then the same sequence extended by one token. For every layer, the keys and
values at positions $1$ to $t$ are identical in the two runs.
\end{proposition}

\begin{proof}
The proof is by induction on the layers. The statement carried from one
layer to the next is that the input vectors of the layer at positions $1$ to
$t$ are identical in the two runs.

\emph{First layer.} The input vectors at positions $1$ to $t$ are the
embeddings of the same tokens at the same positions in both runs, so they are
identical.

\emph{From one layer to the next.} Suppose that the input vectors of a layer
at positions $1$ to $t$ are identical in both runs. Keys and values are
computed from each position's own input vector, so they are identical at
those positions, which is the claim for this layer. Attention at a position
$i \le t$ reads only keys and values at positions up to $i$, because of the
mask, so its result is identical as well, and the remaining steps of the
layer act on each position separately. Hence the output vectors at positions
$1$ to $t$, which are the input vectors of the next layer, are identical, and
the induction continues.
\end{proof}

Two consequences follow. Reusing the cache gives the same output as
recomputing it, up to the rounding of the arithmetic, so it is an
optimization with no loss. Moreover, the cache is a function of the prompt and the
weights alone. Any machine with the same weights computes the same cache,
and a cache computed on one machine can be used on another. The second consequence is what makes the split studied here possible.

The cache is large. Its size is counted in the next chapter: for the model used here, 80 KiB per token.

\subsection{Two Phases with Different Bottlenecks}

Inference has two phases. In \emph{prefill} the model processes the whole
prompt at once and fills the cache. All positions are computed in parallel,
and the GPU's arithmetic units are the limit. In \emph{decode} the model
produces the answer one token at a time. Each step computes very little and
reads the whole cache, so the limit is how fast the GPU can read its own
memory \citep{pope2023}.

The two phases place different demands on the hardware and interfere with each other when they share it: a long prefill delays the decode steps of
every other request on the same GPU. Running them on separate machines, each
tuned for its phase, improves the service that a given amount of hardware can
deliver \citep{patel2024splitwise,zhong2024distserve}.

\subsection{Cache Blocks}

A serving engine handles many requests of unpredictable length, and giving
each a contiguous piece of GPU memory for its cache would waste most of it.
The engine used here divides cache memory into fixed blocks of 16
tokens and gives each request a list of blocks, as an operating system gives
a process pages \citep{kwon2023vllm}. Whatever moves the cache between
machines therefore moves blocks, not one tensor.

\subsection{User-Visible Latency}

Two latencies describe the experience of a user of a language model. \emph{Time
to first token} is the delay from sending the prompt to seeing the first
piece of the answer. It contains the whole prefill. \emph{Inter-token latency}
is the gap between later tokens, and it is set by decode. A system can be
good at one and poor at the other, so both are reported.

\subsection{Disaggregating the Two Phases}

In disaggregated serving, a prefill server computes the cache for a prompt
and a decode server generates the answer from it. Between the two, the cache
has to cross a network. Several systems are built on this pattern
\citep{zhong2024distserve,patel2024splitwise,qin2024mooncake,dynamo}, and
the serving engine used here supports it through a connector interface with
a transfer library underneath \citep{vllm_disagg,nixl}.

The split adds one term to the time to first token: the transfer
(\cref{eq:ttft}). For a 2,048-token prompt on the model used here that is 160 MiB.
The next chapter computes what those bytes cost on commodity Ethernet,
examines what else crosses the link with them, and shows why the order in
which they are sent matters.
